\documentclass[10pt,a4paper]{amsart}
\usepackage[margin=19mm]{geometry}
\usepackage{amsmath,amssymb,mathtools}
\usepackage[scr=rsfs,cal=euler]{mathalfa}
\usepackage{xfrac}
\usepackage[unicode,hidelinks,bookmarks=false]{hyperref}
\newcommand{\ie}{i.e.\ }
\newcommand{\cf}{cf.\ }
\newcommand{\resp}{respectively\ }
\newcommand{\Subsection}{\subsection}
\providecommand{\nobreakdash}{\nobreak\hbox{-}}
\setlength{\emergencystretch}{2em}
\multlinegap0pt
\usepackage{bm}
\usepackage{bbm}
\usepackage{dsfont}
\usepackage{xspace}
\usepackage{cancel}
\usepackage{mathtools}
\usepackage{cleveref}
\usepackage{amsthm}
\makeatletter
\@ifundefined{thm}{\newtheorem{thm}{Theorem}}{}
\@ifundefined{assumption}{\newtheorem{assumption}[thm]{Assumption}}{}
\@ifundefined{prop}{\newtheorem{prop}[thm]{Proposition}}{}
\makeatother
\DeclareMathOperator{\diam}{diam}
\title[Quantifying the noise sensitivity of the Wasserstein metric ]{Quantifying the noise sensitivity of the Wasserstein metric for images}
\author{Erik Lager, Gilles Mordant, Amit Moscovich}
\date{Source: arXiv:2510.01015v3 (CC BY 4.0)}
\begin{document}
\maketitle

\noindent\emph{Source and licence.} Quantifying the noise sensitivity of the Wasserstein metric for images. Version arXiv:2510.01015v3 (CC BY 4.0), distributed under the Creative Commons Attribution 4.0 International licence, as recorded in the arXiv licence metadata for this exact version and shown on the abstract page https://arxiv.org/abs/2510.01015v3. Full rights notice: licenses/arxiv-2510.01015-CC-BY-4.0.txt.\par\smallskip

\noindent\textit{Editorial scope.} Counted unit: the Thm whose source label is \texttt{thm: KR between a measure and its noisy counterpart} in the article (source0.tex lines 1450--1461, source label \texttt{thm: KR between a measure and its noisy counterpart}), delivered together with the article's own complete original proof of it (source0.tex lines 1463--1505, one \texttt{proof} environment of 43 lines). No mathematics is written, repaired or re-derived: statement and proof are the article's own text line for line. Required context retained uncounted: assum: Noise (lines 300--310); prop: KR (lines 1417--1438). Every definition, display and label of those lines is retained unchanged. Omitted from this package and not redistributed: 1--299, 311--1416, 1439--1449, 1462--1462, 1506--1506. Nothing inside a retained excerpt is omitted or altered: every retained body is delivered exactly as the article's lines are written, so no omission receipt of any kind accompanies this package. Citations: the retained excerpts carry no citation command at all, so no bibliography entry of the article is transcribed and no reference is redistributed. Reference adaptation: none was needed and none was made; every reference inside the delivered unit resolves inside it, so no unresolved reference or citation is produced. Printed slips kept literal: no printed word, symbol, hypothesis or number of the counted statement or of its proof was altered. Layout and class: the article's own class is \texttt{article} with packages including microtype, graphicx, subcaption, booktabs, hyperref, icml2026, amsmath, amssymb, mathtools, amsthm, url, soul, doi, placeins; the wrapper is the lane's pinned \texttt{amsart} editorial wrapper at 10pt on A4, which restates the theorem-like environments the retained text needs and adds the article's own macro definitions that the retained text needs (\texttt{\textbackslash diam}), with extra wrapper packages (bm, bbm, dsfont, xspace, cancel, mathtools, cleveref). The delivered numbering is the unit's own. Assets: no retained span contains a figure environment, a graphics-inclusion command, a PDF-page inclusion, a table or a TikZ picture; no image file is copied into this package, the package declares no asset dependency and no image file is redistributed with it. Licence: version arXiv:2510.01015v3 of this article is distributed under the Creative Commons Attribution 4.0 International licence, as recorded in the arXiv licence metadata for this exact version and shown on the abstract page https://arxiv.org/abs/2510.01015v3. A full rights notice is delivered at licenses/arxiv-2510.01015-CC-BY-4.0.txt. Characters: every retained excerpt is pure ASCII, so the delivered engine, XeLaTeX with its default Latin Modern text font, reports no missing character.
\begin{assumption}
\label{assum: Noise}
Consider an image modeled as an $n \times n$ grid of pixels and set $m= n^2$. Assume that the noise vector $N = (N_1, \ldots, N_m)$ is drawn from a multivariate normal distribution $\mathcal{N}(0, \Sigma)$, where the covariance matrix $\Sigma$ is an $m \times m$ matrix defined as
\begin{align}
\Sigma_{ij} =
\begin{cases}
    \sigma^2 & \text{if } i=j \\
    -\frac{\sigma^2}{m-1} & \text{if } i \ne j.
\end{cases}
\end{align}
\end{assumption}

\begin{prop}
\label{prop: KR}
Consider a square $2^\eta\times 2^\eta$ grid, where $\eta \ge 0$ is integer, and a dyadic partition scheme. Let $\mu, \nu$ be two measures on the grid, not necessarily with equal masses. 
It holds that,
\begin{align}
&\operatorname{KR}_{p,C}(\mu, \nu)  \le
        \frac{C^p}{2}
            \lvert m_\mu - m_\nu\rvert
            \\&\qquad+
           \diam(S)^p  2^{3p-1}
            \sum_{k=\ell^*}^{\eta}
            2^{-pk}
            \sum _{Q\in \mathcal{D}^k}
            \lvert
                \mu(Q)-\nu(Q)
            \rvert.    \nonumber
\end{align}
where
\(
\ell^*=1+ \min\left( L, \left\lfloor \max\big(0 , \log_2 \left( 2\diam(S)/C  \right)\big) \right\rfloor \right).
\)
\end{prop}

\begin{thm}
\label{thm: KR between a measure and its noisy counterpart}
Let $\mu : G_n \to [0,1]$ be a probability measure on the $n\times n$ unit grid $G_n$ with cyclic boundary conditions, and let $\varepsilon$ satisfy Assumption~\ref{assum: Noise}. Further assume that $n = 2 ^\eta$, for $\eta \in \mathbb{N}$. Then, for $C>0$,
\begin{align}
\label{eq: KR_cases}
    &\mathbb{E} \operatorname{KR}_{p,C}^\pm(\mu+\varepsilon,\mu) \\
    &\le \begin{cases} 
          \diam(S)^p 2^{3p-1} \frac{\sqrt{2}\sigma}{\sqrt{\pi}} n (\eta-\ell^*+1) & \text{if } p=1, \\[1ex]
        \diam(S)^p 2^{3p-1} \frac{\sigma}{\sqrt{2\pi}}  n (2^{2-\ell^*} - 2^{1-\eta}) & \text{if } p=2.
    \end{cases} \nonumber
\end{align}
\end{thm}

\begin{proof}[Proof of \cref{thm: KR between a measure and its noisy counterpart}]



We apply Proposition~\ref{prop: KR} to the probability measures $\mu_-+(\mu+\varepsilon)_+$ as well as $(\mu+\varepsilon)_-+ \mu_+$.  
First, note that 
\begin{align}
    &\mathbb{E} \lvert m_{ \mu_-+(\mu + \varepsilon)_+} - m_{\mu_+ +  (\mu + \varepsilon)_-} \rvert
    \\
    &=\mathbb{E} \bigg\lvert \sum_{x} (\mu_- - \mu_+)(x) + (\mu + \varepsilon)_+(x) - (\mu + \varepsilon)_-(x)\bigg\rvert = 0. \nonumber
\end{align}
Taking expectations and using independence and zero mean of the noise,
\begin{align}
    &\sum_{k=\ell^*}^{\eta}
            2^{-pk}
            \sum _{Q\in \mathcal{D}^k}
            \lvert
                \mu_-(Q)+ (\mu + \varepsilon)_+(Q)
                \\&-(\mu + \varepsilon)_-(Q) - \mu_+(Q)
            \rvert  
            = \sum_{k=\ell^*}^{\eta}
            2^{-pk}
            \sum _{Q\in \mathcal{D}^k}
            \lvert
                 \varepsilon(Q)
            \rvert. \nonumber
\end{align}

Since each $\varepsilon(x)$ is Gaussian with variance $\sigma^2$, one has $\mathbb{E}|\sum_{x\in Q}\varepsilon(x)| \le \sigma \sqrt{|Q|}\sqrt{2/\pi}$. Furthermore, the dyadic family $\mathcal{D}_k$ has $|\mathcal{D}_k|=2^{2k}$ cubes of cardinality $|Q|=n^2/2^{2k}$. Therefore
\begin{align}
    &\sum_{Q\in\mathcal{D}_k} \mathbb{E}\Big| \sum_{x\in Q} \varepsilon(x) \Big|
    \le \sigma \sqrt{\tfrac{2}{\pi}} \sum_{Q\in\mathcal{D}_k} \sqrt{|Q|}
    \\&= \sigma \sqrt{\tfrac{2}{\pi}} \cdot 2^{2k} \cdot \frac{n}{2^{k}}
    = \sigma \sqrt{\tfrac{2}{\pi}}\, n\, 2^{k}. \nonumber
\end{align}
Plugging this into the multiscale sum yields
\begin{equation}
 \sum_{k=\ell^*}^{\eta} 2^{-pk} \sum_{Q\in\mathcal{D}_k} \mathbb{E}\Big| \sum_{x\in Q} \varepsilon(x) \Big|
    \le  \sigma \sqrt{\tfrac{2}{\pi}}\, n \sum_{k=\ell^*}^{\eta} 2^{(1-p)k}. \qedhere
\end{equation}


\end{proof}

\end{document}
